The first line contains a single integer $t$ --- the number of test cases.
$(1 \le t \le 10)$

Each test case consists of the following:

The first line contains three integers $n$, $m$, and $q$ --- the number of nodes, edges, and queries, respectively.
$(2 \le n  \le 1000,\ 1 \le m \le 1000,\ 1 \le q \le 10^4)$

The next $m$ lines each contain three integers $u$, $v$, and $w$ --- representing an undirected edge between nodes $u$ and $v$ with weight $w$.
$(1 \le u, v \le n,\ 1 \le w \le 10^9)$

The next line contains a string of length $n$ consisting only of characters `.' and `x', representing the initial status of each location (`.' indicates that the location is open) and (`x' indicates that the location is blocked).

The next $q$ lines each contain a query starting with $a_i\ $--- where:

\begin{itemize}
\item If $a_i = 0$, the format of query is $a_i\ x_i\ y_i$ --- it is a path query: compute the shortest path from node $x_i$ to node $y_i$ using only unblocked nodes.

\item If $a_i = 1$, the format of query is $a_i\ x_i\ $ --- it is a toggle operation: toggle the blocked/unblocked status of node $x_i$. 
\end{itemize}

\textbf{Note}:
\begin{itemize}
\item The sum of $n \times m$ over all test cases does not exceed $10^6$.

\item For each test case, the total sum of $a_i$ over all queries is at most $10$.

\item In the string `.' is the dot or fullstop in english keyboard and `x' is the alphabetic character x.
\end{itemize}